\section{Experimental Setup}
We use  \texttt{Llama-3.1-8B-Instruct}, \texttt{google/gemma-7b-it}, and \\ \texttt{mistralai/Mistral-7B-Instruct-v0.3} as they are open-source instruction tuned models~\cite{grattafiori2024llama3,gemmateam2024gemmaopenmodelsbased}. We opt to use smaller models for broader accessibility and lower computational costs. We extract the control vectors for layers $16,17,18$ of the model, as middle layers were shown to have the best effect on the downstream output \cite{rimsky-etal-2024-steering}. We use \{-1, -0.5, 0,  0.5, 1\} as the set of strengths and set the temperature to 0.0. In preliminary experiments, we found that using a control vector strength greater than $|1|$ caused degeneration. For the main results reported in~\Cref{tab:main-results-llama}, ~\Cref{tab:main-results-gemma} and ~\Cref{tab:main-results-mistral}, we randomly sample 1024 examples from the PLOS and SciNews test sets for evaluation. We use the full eLife and Eureka test datasets for evaluation, as they contain less than 1024 examples (241 and 1010 respectively). We ran all of our experiments on a grid with NVIDIA A100 and H200 GPUs available. Each experiment only requires a single A100 GPU.

Past work has shown that the Coleman-Liau readability index (CLI)~\cite{Coleman1975ACR} and Dale-Chall readability scores (DCRS)~\cite{dale1948formula} correlate the highest with human judgments of readability~\cite{cachola-etal-2025-evaluating}. To avoid overfitting to our evaluation, we use DCRS for our data selection (\S\ref{sec:data-selection}) and CLI for our system evaluation (\S\ref{sec:4-chapter-results}). To evaluate system performance, we report the Pearson~\cite{PearsonVIINO} and Kendall-Tau~\cite{Kendall1938ANM} correlation of the specific strength scalar with CLI. CLI provides a lower score for higher readability, while our setup treats the lay summary as the positive example. In other words, we expect higher control vector strengths to produce lower CLI scores. Therefore, we multiply the CLI scores by $-1$ so the resulting correlations are positive. We also use BERTScore \cite{zhang2020bertscore} with the abstract as the reference, since some dataset lay summaries contain extraneous information not present in the original document and thus make poor evaluation references.

% \subsection{Data}
% The eLife dataset consists of 5,000 full-text biomedical research articles paired with non-technical lay summaries published by the eLife scientific journal. The Eureka dataset is derived from the EurekAlert corpus provided by the HTSS project. It contains 5,000 general scientific articles along with corresponding simplified lay summaries written as news reports for the public. The PLoS dataset comprises of around 40,000 full-length biomedical research articles paired with expert-written lay summaries, sourced from various journals published by the Public Library of Science (PLOS). The SciNews dataset includes over 40,000 scientific articles from nine distinct domains paired with a corresponding news report written for a lay audience. 

\subsection{Baselines Implementation Details}\label{sec:base-details}
We compare the ability of the control vector methods to In-Context Learning (ICL), Supervised Fine Tuning (SFT), LoRa/QLoRa SFT and DSPy Prompt Optimization. We use the sub-selected 128 most readable summaries dataset to ensure data fairness with our method.
 ICL has shown great abilities in instruction following, and is a similarly low data, lightweight method~\cite{Brown2020LanguageMA}. For ICL, we take a positive and negative example from the extraction set (\S\ref{sec:data-selection}) as examples of a lay summary and technical summary, respectively. We then provide the document to the model and ask it to summarize the document with a specified readability level. We repeat this with 5 levels of specified readability. We additionally use SFT and QLoRa/LoRa SFT as baselines \cite{hu2021loralowrankadaptationlarge, dettmers2023qloraefficientfinetuningquantized}. Prompt Optimization has shown promise as a way to control language models without expensive training \cite{khattab2023dspycompilingdeclarativelanguage}. 
We compared our control-vector method against commonly used alternatives, including In Context Learning (ICL) \cite{Brown2020LanguageMA}, Supervised Fine-Tuning (SFT), LoRA/QLoRA based Supervised Fine-Tuning \cite{hu2021loralowrankadaptationlarge, dettmers2023qloraefficientfinetuningquantized}, and Prompt Optimization.
\subsubsection{ICL}
We use the following prompt for our ICL baseline, based on the prompts used in~\cite{ribeiro-etal-2023-generating}:

\textexample{
Summary 0 properties: 0\% readability \textbackslash n \\
Summary 0: <TECHNICAL SUMMARY> \textbackslash n \textbackslash n \\
Summary 1 properties: 100\% readability \textbackslash n \textbackslash n \\
Summary 1: <LAY SUMMARY> 
\textbackslash n \textbackslash n \\
Summarize the input document with the following properties: <0,25,50,75,100>\% readability \textbackslash n \\
Input document: <SOURCE> \textbackslash n \\
Summary:}

We use the following prompt during extraction of the control vectors:

\textexample{Document: <SOURCE>  \textbackslash n 
Summarize: <TARGET>}

We use the following prompt during inference using the control vectors:

\textexample{Document: <SOURCE>  \textbackslash n 
Summarize:}

We experiment with the effect of number of examples provided to the model and its effect on controllability. In our setup, a ``shot'' is a pair of positive and negative summaries, so for each setting a $N$-shot experiment contains $2*N$ total summaries. We experiment with $N=[1,8,64,128]$ shots to match up to the total number of example pairs provided to the control vector experiments. The shots are randomly sampled from the training dataset and we randomly sample 100 examples from the test set for inference. We use LLaMa 8b, the best performing model in~\Cref{tab:main-results}. The results are in~\Cref{tab:icl-mini}.
% Please add the following required packages to your document preamble:
% \usepackage{subcaption}
\begin{table}[h!]
    \centering
    \scriptsize
\begin{subtable}[t]{0.48\textwidth}
\centering
\begin{tabular}{l|lll}
\textbf{N} & \textbf{Pearson} & \textbf{Kendall} & \textbf{BertS F1} \\ \hline
1 & -0.066 & -0.039 & 0.845 \\
8 & -0.023 & -0.013 & 0.841 \\
64 & 0.016 & 0.006 & 0.838 \\
128 & 0.027 & 0.022 & 0.841 \\
\end{tabular}
\caption{Elife}
\label{tab:icl-mini-elife}
\end{subtable}
\hfill
\begin{subtable}[t]{0.48\textwidth}
\centering
\begin{tabular}{l|lll}
\textbf{N} & \textbf{Pearson} & \textbf{Kendall} & \textbf{BertS F1} \\ \hline
1 & -0.062 & -0.035 & 0.836 \\
8 & -0.018 & -0.017 & 0.831 \\
64 & -0.001 & 0.001 & 0.834 \\
128 & 0.038 & 0.035 & 0.844 \\
\end{tabular}
\caption{Eureka}
\label{tab:icl-mini-eureka}
\end{subtable}

\vspace{1em}

\begin{subtable}[t]{0.48\textwidth}
\centering
\begin{tabular}{l|lll}
\textbf{N} & \textbf{Pearson} & \textbf{Kendall} & \textbf{BertS F1} \\ \hline
1 & -0.030 & -0.021 & 0.847 \\
8 & -0.023 & -0.014 & 0.838 \\
64 & 0.011 & 0.004 & 0.837 \\
128 & -0.019 & -0.007 & 0.843 \\
\end{tabular}
\caption{PLOS}
\label{tab:icl-mini-plos}
\end{subtable}
\hfill
\begin{subtable}[t]{0.48\textwidth}
\centering
\begin{tabular}{l|lll}
\textbf{N} & \textbf{Pearson} & \textbf{Kendall} & \textbf{BertS F1} \\ \hline
1 & -0.028 & -0.024 & 0.835 \\
8 & -0.009 & -0.009 & 0.831 \\
64 & -0.001 & -0.011 & 0.830 \\
128 & -0.013 & -0.007 & 0.834 \\
\end{tabular}
\caption{SciNews}
\label{tab:icl-mini-scinews}
\end{subtable}
\caption[Results using In-Context Learning, varying the number of shots.]{Results using In-Context Learning, varying the number of shots. Additional shots do not generally increase performance.}
\label{tab:icl-mini}
\end{table}

The number of shots provided for In-Context Learning does not have a significant effect on the overall controllability or quality of the summaries. Therefore, to save compute by number of tokens, we use $N=1$ for the ICL experiments in the main body of results, reported in~\Cref{sec:4-chapter-results}.


\subsubsection{SFT}
To train our model using supervised fine-tuning (SFT), we leveraged the datasets originally generated for learning control vectors. For each source document, we included both the positive and negative examples as separate training instances paired with the same reference summary, effectively doubling the size of the training set.

\textexample{Document: <SOURCE>\textbackslash n  Summarize with readability level <READABILITY>.\textbackslash n  Summary: <SUMMARY>\textbackslash n \textbackslash n}\label{baseline_prompt}

We used the prompt shown above for training. For each example, we computed the Dale--Chall readability score of the reference summary and used this value as the target readability level specified in the prompt. The model was trained using a causal language modeling loss applied only to the summary portion of the output, implemented using the TRL framework \cite{vonwerra2022trl}. During training, we did not apply output truncation and instead relied on the maximum input length supported by the base model.

\subsubsection{LoRA/QLoRA SFT}
Our second baseline uses LoRA and QLoRA adapters for parameter-efficient fine-tuning. The training procedure is identical to SFT, except that only adapter parameters are updated. We primarily use LoRA; when memory constraints prevent LoRA training, we instead employ 4-bit QLoRA to enable training under reduced memory usage.

\subsubsection{Prompt Optimization}
For prompt optimization, we used the DSPy \cite{khattab2023dspycompilingdeclarativelanguage} framework to automatically select an optimized prompt. We provided DSPy with the same training examples and initial prompt used for SFT. The optimization objective was the mean squared error between the Dale--Chall readability score of the generated summary and the target readability specified in the prompt. We performed optimization using the MiPROv2 \cite{opsahlong2024optimizinginstructionsdemonstrationsmultistage} optimizer with default hyperparameters.


% \subsection{Models}
% We used \texttt{google/gemma-7b-it} \cite{gemmateam2024gemmaopenmodelsbased} and \texttt{meta-llama/Llama-3.1-8B-Instruct} \cite{grattafiori2024llama3} as they are strong open source models instruction following models that are easily trainable on 1-2 A100 80GB GPUs. 

\subsubsection{Hyperparameters Used}
We took hyperparameters from the AxBench Paper \cite{wu2025axbenchsteeringllmssimple} for the LoRa-SFT and SFT as they provided a fair comparison between Control Vectors, SFT, and Prompting-based Methods. However, we modified the number of epochs and batch size as our datasets were much longer in both number of items and length of inputs. In Table \ref{tab:baseline_hyperparameters}, we show the hyperparameters used for each baseline method. 

\begin{table*}[t]
\centering
\scriptsize
\begin{tabular}{lcccc}
\toprule
\diagbox{Hyperparameter}{Baseline Method} & LoRa & QLoRa & SFT & DSPy \\
\midrule
Learning rate              & 5.00E-03 & 5.00E-03 & 4.00E-05 & -- \\
Epochs                     & 3        & 3        & 3        & -- \\
Quantization               & --    & 4-bit    & --       & -- \\
Batch size                 & 1        & 1        & 1        & -- \\
LoRA rank                  & 4        & 4        & --       & -- \\
LoRA alpha                 & 32       & 32       & --       & -- \\
LoRA dropout               & 0.1      & 0.1      & --       & -- \\
Target modules             & o\_proj  & o\_proj  & --       & -- \\
Layers to transform        & 12, 20, 31, 39 & 12, 20, 31, 39 & -- & -- \\
Optimizer                  & AdamW       & AdamW       & AdamW       & MIPROv2 \\
Max bootstrapped demos     & --       & --       & --       & 4 \\
Max labeled demos          & --       & --       & --       & 4 \\
Auto                        & --       & --       & --       & medium \\
\bottomrule
\end{tabular}
\caption{Hyperparameters for each baseline method.}
\label{tab:baseline_hyperparameters}
\end{table*}

\subsubsection{Inference and Evaluation}
For SFT, QLoRa/LoRa SFT and DSPy Prompt Optimization, we tested the model's ability by telling the model to provide summaries at readability scores of 1,3,5,7,9 mirroring the range of Dale Chall, to evaluate the model's ability to change the readability similarly to the control vectors. 





\section{Meeting Data Requirements}\label{sec:data-selection}

We begin with four different scientific summarization datasets: eLife~\cite{goldsack2022making}, Eureka~\cite{Zaman2020HTSSAN}, PLOS~\cite{goldsack2022making}, and SciNews~\cite{pu-etal-2024-scinews}. These datasets are designed for  scientific PLS, allowing us to use the lay summary as the positive example and the abstract as the technical, negative example. 

% meet 2 requirements: (1) The positive and negative examples must be sufficiently different in readability and (2) The positive and negative examples must have approximately the same content, only differing in their readability level.


We choose an extraction dataset that follows the requirements outlined in~\S\ref{sec:4-chapter-method}. Req. 1 states that the paired examples must be sufficiently different in readability. We plot the Dale-Chall scores of the positive and negative examples for each dataset in a histogram. We additionally calculate the Bhattacharyya distance, which measures the overlap between two distributions~\cite{Bhattacharyya1943OnAM}. A dataset that best meets this requirement will have a large visual separation between the two distributions, as well as a high Bhattacharyya distance. \Cref{fig:dale-chall-hists} contains both the histograms and the Bhattacharyya distance for each dataset. eLife has both the highest visual separation and highest Bhattacharyya distance. PLOS has a particularly low separation, indicating that the PLOS lay summaries are not significantly more readable than the technical summaries.
\begin{figure}[!ht]
    \vspace{-3mm}
    \centering
    \includegraphics[width=.75\linewidth]{research-chapters/4-chapter/figures/dale_chall_histograms.png}
    \vspace{-3mm}
    \caption[Histogram of the Dale-Chall readability scores for the lay and technical summaries.]{Histogram of the Dale-Chall readability scores for the lay and technical summaries and the Bhattacharyya distance (BD) of each dataset.}
    \label{fig:dale-chall-hists}
    \vspace{-5mm}
\end{figure}

Req. 2 states that the positive and negative examples must have approximately the same content, only differing in their readability level. In order to measure overlap in content, we use Sentence Transformers to retrieve an embedding matrix for the example pairs, then compute the cosine similarity~\cite{reimers-gurevych-2019-sentence}. We use Specter for the embeddings, a model trained on scientific data~\cite{cohan-etal-2020-specter}. The resulting mean cosine similarity scores are as follows: $0.631$ for eLife, $0.604$ for PLOS, $0.575$ for Eureka, and $0.590$ for SciNews. Based on these scores, we conclude that eLife and PLOS have the highest content similarity while Eureka and SciNews have the lowest. This aligns with the construction of the datasets: eLife and PLOS's lay summaries are written by the original paper authors to be accessible explanations of their papers while Eureka and SciNews's summaries are news reports discussing papers. Upon inspection, we find that Eureka and SciNews often contain information not present in the abstract, such as interviews with the authors or information on a study's funding source. 
Based on this analysis, we conclude that eLife best meets the requirements. As the control vector method does not require a large extraction dataset, we select 128 of the most readable examples (lowest Dale-Chall) from the eLife training set. Choosing the most readable examples gives the control vectors the best examples of readability. We use this setup to report our main results in~\S\ref{sec:4-chapter-results}. 

\input{research-chapters/4-chapter/tables/dataset-examples}


We include examples of each dataset in~\Cref{tab:dataset-examples-elife} (eLife), \Cref{tab:dataset-examples-plos} (PLOS), \Cref{tab:dataset-examples-eureka} (Eureka), and \Cref{tab:dataset-examples-scinews}. In the eLife example, the lay summary is significantly more readable than the technical summary, using less technical language, while only focusing on the content of the paper. This is an example of a pair of summaries that are compliant with the requirements outlined in~\S\ref{sec:4-chapter-method}. In the PLOS example, although it is clear that both summaries cover the same material, the lay summary is not significantly more readable than the technical summary, including highly complex terms such as ``glycosaminoglycans.'' This example meets Req. 2 but not Req. 1. The Eureka and SciNews examples contain extraneous information not present in the technical summary. Both example lay summaries contain interviews with the papers' authors, meaning these summaries do not meet Req. 2. We note that although this extraneous information is problematic for control vectors, it is also likely problematic for any method. At best, the extraneous information will make a method ineffective. At worst, it could teach a model to hallucinate interviews that did not happen. For this reason, we urge future researchers to use caution using these datasets for summarization training.

